\begin{Problem}[Density Matrix of a Single Qubit]
For a given ensemble of qubits the observables
\[
A=
\begin{pmatrix}
3 & 0\\
0 & 1
\end{pmatrix},
\qquad
B=
\begin{pmatrix}
1 & 1\\
1 & -1
\end{pmatrix},
\qquad
C=
\begin{pmatrix}
0 & 2i\\
-2i & 0
\end{pmatrix},
\]
have been measured many times, giving the expectation values
\[
\langle A\rangle = 2,
\qquad
\langle B\rangle = \frac12,
\qquad
\langle C\rangle = 0.
\]

\begin{enumerate}
\item[(a)] Determine the density matrix of the ensemble.
\item[(b)] Is the ensemble pure or mixed?
\item[(c)] If we measure $\sigma^z$, what is the probability to get $+1$?
\item[(d)] Determine the expectation values
\[
\langle \sigma^x\rangle,\qquad
\langle \sigma^y\rangle,\qquad
\langle \sigma^z\rangle.
\]
\end{enumerate}
\end{Problem}
\begin{proof}
	\begin{parts}
	\item Given a density matrix $\rho$, the expectation value of an operator $\mathbf{A}$ is given by $\langle \mathbf{A}\rangle = \Tr(\rho \mathbf{A})$. We implement the trace 1 and hermiticity conditions by defining our density matrix using
		\[
			\rho = \begin{pmatrix} a & b \\ b^* & 1-a \end{pmatrix} 
		.\] 
		Substituting this in, we find
		\begin{align*}
			\langle \mathbf{A}\rangle &= 1+2a = 2 \\
			\langle \mathbf{B}\rangle &= -1+2a+b+b^* = \frac{1}{2} \\
			\langle \mathbf{C}\rangle &= -2ib + 2ib^*=0 
		\end{align*}
		From the third equation, we see that $b=b^*$, i.e. $b$ is real. From the first equation, we can see easily that $a = \frac{1}{2}$. Substituting this in to the second equation, we find that $2b = \frac{1}{2}$, or $b = \frac{1}{4}$.
	\item We do this using the trace criterion: We have
		\[
			\rho^2 = \begin{pmatrix} \frac{5}{16} & \frac{1}{4} \\ \frac{1}{4} & \frac{5}{16} \end{pmatrix} 
		.\] 
		Thus, the trace is
		\[
		\Tr(\rho^2) = \frac{10}{16}\neq 1
		.\] 
		The ensemble is mixed.
	\item ..
	\item We compute
		\[
			\rho \sigma^z = \begin{pmatrix} \frac{1}{2} & -\frac{1}{4} \\ \frac{1}{4} & -\frac{1}{2} \end{pmatrix} 
		.\] 
		Thus, we have
		\[
		\langle \rho\sigma^z\rangle = 0
		.\] 
		Similarly, we have
		\[
			\rho\sigma^x = \begin{pmatrix} \frac{1}{4} & \frac{1}{2} \\ \frac{1}{2} & \frac{1}{4} \end{pmatrix} 
		\]
		with
		\[
			\langle \rho \sigma^x\rangle = \frac{1}{2}
		\]
		and
		\[
			\rho \sigma^y = \begin{pmatrix} \frac{i}{4} & -\frac{i}{2} \\ \frac{i}{2} & -\frac{i}{4} \end{pmatrix} 
		\]
		with
		\[
		\langle \rho\sigma^y \rangle = 0
		.\] 
	\end{parts}
\end{proof}

\begin{Problem}[Criterion for the Equivalence of Ensembles]
Consider a statistical ensemble $\{|\psi_i\rangle,p_i\}$ with the associated
density matrix
\[
\rho=\sum_{i=1}^{N} p_i |\psi_i\rangle\langle\psi_i|.
\]

Here $N$ is arbitrary (smaller or larger than the dimension of the Hilbert
space) and the vectors $|\psi_i\rangle$ are pairwise different and normalized
but not necessarily orthogonal. Moreover, let
\[
\rho=\sum_{k=1}^{K}\lambda_k |k\rangle\langle k|
\]
be the spectral decomposition of $\rho$ with orthonormal eigenvectors $|k\rangle$
and eigenvalues $\lambda_k$.

The aim of this exercise is to work out an equivalence theorem for different
ensembles with the same density matrix.

\begin{enumerate}
\item[(a)] Show that every vector $|\psi_i\rangle$ can be written as a linear
combination of the eigenvectors $|k\rangle$, i.e.
\[
|\psi_i\rangle\in \mathrm{span}\{|k\rangle\}.
\]

\item[(b)] Use (a) to prove that every $|\psi_i\rangle$ can be written as
\[
\sqrt{p_i}\,|\psi_i\rangle
=
\sum_{k=1}^{K}
c_{ik}\sqrt{\lambda_k}\,|k\rangle,
\]
with coefficients $c_{ik}$ obeying
\[
\sum_{i=1}^{N} c_{ik}c_{il}^{*}
=
\delta_{kl},
\qquad
k,l=1,\dots,K.
\]

\item[(c)] Because of (a) we know that $N\ge K$. For $N=K$ the result of (b)
would imply that the coefficient matrix $c_{ik}$ is unitary. Show that for
$N>K$ the rectangular matrix $c_{ik}$ can be extended to a square unitary
matrix such that the relation in (b) still holds.

\item[(d)] So far we have proven that if
\[
\sum_{i=1}^{N} p_i |\psi_i\rangle\langle\psi_i|
=
\sum_{k=1}^{K}\lambda_k |k\rangle\langle k|,
\]
then there exists a unitary $N\times N$ matrix $c_{ik}$ such that
\[
\sqrt{p_i}\,|\psi_i\rangle
=
\sum_{k=1}^{N}
c_{ik}\sqrt{\lambda_k}\,|k\rangle.
\]
Prove the converse statement.

\item[(e)] Use (a)--(d) to show that
\[
p_i=\sum_{k=1}^{N} T_{ik}\lambda_k,
\]
(or in vector form $\vec p=T\vec\lambda$),
where $T$ is a doubly-stochastic matrix, i.e. all matrix entries are real
and nonnegative and
\[
\sum_{i=1}^{N}T_{ik}
=
\sum_{k=1}^{N}T_{ik}
=
1.
\]

\item[(f)] Use e.g. Mathematica to work out the explicit example
\[
p_1=p_2=\frac12,
\qquad
|\psi_1\rangle
=
\frac1{\sqrt2}
\begin{pmatrix}
1\\
1
\end{pmatrix},
\qquad
|\psi_2\rangle
=
\begin{pmatrix}
0\\
1
\end{pmatrix},
\]
with $N=K=2$, and determine
$\lambda_{1,2}$, $|\phi_{1,2}\rangle$, $c_{ij}$, and $T_{ij}$.
Check the properties derived in (e).
\end{enumerate}
\end{Problem}
\begin{proof}
	\begin{parts}
	\item This follows from the spectral theorem; since $\rho$ is Hermitian, the eigenvectors $\ket{k}$ form an orthogonal decomposition of the Hilbert space.
	\end{parts}
\end{proof}
